% week10_tue_probability.tex -- DSCT Week 10 Tuesday, Oct 20 2026.
% Sources: week-10/instructor/tuesday-run-of-show.md, week-10/README.md,
% week-10/student/tuesday-activity.md, lessons/08-probability.md.
% Every number below was recomputed in Python (seed 10 for the simulations);
% see the verification notes in the speaker notes.
\documentclass[aspectratio=169,10pt]{beamer}
\usetheme{Madrid}
\usepackage{listings}
\lstset{basicstyle=\ttfamily\small,breaklines=true,columns=fullflexible,frame=single}
\title{A Positive Result Is Not a Probability of Truth}
\subtitle{Week 10 Tuesday --- Probability, uncertainty, and evidence}
\author{COMSC 2043 \textbullet\ Discrete Structures and Critical Thinking}
\date{October 20, 2026}
\begin{document}

\frame{\titlepage
\note{75-minute meeting. Spine: Sources, Rules/Assumptions, Work, Check, One-Sentence Summary. Do not let a gambling example turn into an unstated real-world frequency claim. Keep P(+ | C) distinct from P(C | +) all hour. Manual tables are enough; the script is a check, not the lesson.}}

\begin{frame}{Today's destination}
\begin{itemize}
\item Examine the claim ``90\% accurate means 90\% likely true.''
\item Build a four-cell count table and calculate $P(C\mid +)$.
\item Change the base rate and predict first.
\item See what a simulation checks, and what it cannot establish.
\item Draft a one-sentence conclusion for the Decision Gate.
\end{itemize}
\note{Run of show: 0:00 AI Fluency claim (8 min). 0:08 sample spaces, events, complements, conditional probability (12). 0:20 screening example (18). 0:38 students calculate then change prevalence (14). 0:52 calculation versus simulation (13). 1:05 preview Thursday and Decision Gate, draft conclusion (12). 1:17 inspect the trace, bridge to Week 11 (8).}
\end{frame}

\begin{frame}{AI Fluency: examine one sentence}
\begin{block}{A claim someone might paste into a chat}
``The test is 90\% accurate, so a positive result means you are 90\% likely to have the condition.''
\end{block}
\begin{itemize}
\item Two different probabilities hide inside ``90\%.''
\item Which one is the test's \emph{sensitivity}? Which one would the claim need?
\item Name the ambiguity before doing any arithmetic.
\end{itemize}
\note{Eight minutes. Have pairs name the conditional probability the sentence confuses. Target: P(+ given C) is what 90 percent describes; the claim needs P(C given +). Write both on the board with the given-bar and leave them up.}
\end{frame}

\begin{frame}{Vocabulary we will use}
\begin{description}
\item[Sample space] every outcome we agree is possible.
\item[Event] a set of outcomes, e.g.\ $C$ = ``has the condition,'' $+$ = ``tests positive.''
\item[Complement] ``not $C$'': everything in the sample space outside $C$.
\item[Conditional] $P(A\mid B)$ = share of $B$ that is also $A$; the denominator changes.
\end{description}
\note{Twelve minutes with the labeled event/count table. The key move: conditioning shrinks the sample space to B, so B is the denominator. Students sketch a two-by-two grid before the next slide.}
\end{frame}

\begin{frame}{Sources and rules for the screening example}
\begin{itemize}
\item \textbf{Source:} a clinic screens 10{,}000 people. Prevalence is 1\%. Sensitivity is 90\%. Specificity is 95\%.
\item \textbf{Sensitivity:} $P(+\mid C)=0.90$. \textbf{Specificity:} $P(-\mid \text{not }C)=0.95$.
\item \textbf{Assumption:} these rates apply to this population.
\item Name one reason that assumption could be unsafe.
\end{itemize}
\note{Use only the stated rates. Possible unsafe-assumption answers from the class: the clinic sees symptomatic people, the rates were measured on a different group, repeat testing is not independent. Accept any bounded one and move on.}
\end{frame}

\begin{frame}{Natural frequencies: the four-cell table}
\centering
\begin{tabular}{l|r|r|r}
 & Tests $+$ & Tests $-$ & Total \\ \hline
Has $C$ (1\%) & 90 & 10 & 100 \\
Not $C$ (99\%) & 495 & 9{,}405 & 9{,}900 \\ \hline
Total & 585 & 9{,}415 & 10{,}000
\end{tabular}
\vspace{3mm}
\begin{itemize}
\item $90\%$ of 100 is 90. $5\%$ of 9{,}900 is 495.
\item Check: $90+10+495+9405=10{,}000$.
\end{itemize}
\note{Model it live: Sources, Rules, Work, Check. Fill the C row first (100 people: 90 test positive, 10 negative), then the not-C row (9,900 people: 5 percent false positives is 495, the rest 9,405). Verified: all four cells and both totals recomputed in Python, and the cells sum to 10,000.}
\end{frame}

\begin{frame}{Reading the table: two different conditionals}
\begin{columns}[T]
\column{.48\textwidth}
\textbf{$P(+\mid C)$ --- sensitivity}\\[1mm]
Denominator: everyone with $C$ (100).\\
$90/100 = 90\%$
\column{.48\textwidth}
\textbf{$P(C\mid +)$ --- what a positive means}\\[1mm]
Denominator: everyone who tested positive (585).\\
$90/585 = 2/13 \approx 15.4\%$
\end{columns}
\vspace{4mm}
\begin{alertblock}{Not the same number}
A positive result here is correct about 15 times in 100, not 90.
\end{alertblock}
\note{This is the center of the meeting. Verified: 90/585 = 2/13 = 0.1538..., and Bayes' rule on the stated rates gives the same 0.1538. Also P(+) = 585/10,000 = 5.85 percent. Ask: where did the other 495 positives come from? From the huge not-C group at a small error rate.}
\end{frame}

\begin{frame}{Check, then change the base rate}
\begin{enumerate}
\item Check: cells total 10{,}000, and the positive column (585) is the denominator.
\item \textbf{Predict first:} prevalence rises from 1\% to 20\%. Does $P(C\mid +)$ go up, down, or stay put?
\item Then calculate with the same sensitivity and specificity.
\end{enumerate}
\vspace{2mm}
\centering
\begin{tabular}{l|r|r|r}
 & Tests $+$ & Tests $-$ & Total \\ \hline
Has $C$ (20\%) & 1{,}800 & 200 & 2{,}000 \\
Not $C$ (80\%) & 400 & 7{,}600 & 8{,}000 \\ \hline
Total & 2{,}200 & 7{,}800 & 10{,}000
\end{tabular}
\note{Students work for about fourteen minutes: they calculate the 1 percent case, write the prediction, then rebuild at 20 percent. Verified: 1,800 + 200 + 400 + 7,600 = 10,000, and P(C given +) = 1,800/2,200 = 9/11 = 0.818. Same test, same rates; only the base rate moved. Reveal the table after they commit to a prediction.}
\end{frame}

\begin{frame}{Base-rate neglect}
\begin{itemize}
\item Same test: $P(C\mid +)\approx 15.4\%$ at 1\% prevalence, $\approx 81.8\%$ at 20\%.
\item The test did not change. The population did.
\item Ignoring how common the condition is, and quoting only the test's error rates, is \textbf{base-rate neglect}.
\item A claim about a positive result needs the base rate in the Sources line.
\end{itemize}
\note{Name the fallacy only after they have seen the numbers move. Keep it about the stated example; do not generalize to real clinics or real tests.}
\end{frame}

\begin{frame}{Independence and the gambler's fallacy}
\begin{itemize}
\item Two events are \textbf{independent} if knowing one does not change the probability of the other.
\item Assumption for the exercise: a fair coin, flips independent.
\item Five heads in a row has probability $(1/2)^5 = 1/32 = 0.03125$ \emph{before} you flip.
\item Given five heads already happened, the next flip is still $1/2$: the coin has no memory.
\end{itemize}
\note{Gambling stays a stated process, not a real-world frequency claim: we are told the coin is fair and independent. The contrast with the screening table: there, conditioning on a positive really does change the probability because the events are dependent. Independence is an assumption we state, not a feeling.}
\end{frame}

\begin{frame}[fragile]{What a simulation checks}
\begin{lstlisting}
import random
random.seed(10)
n, pos, pos_c = 1_000_000, 0, 0
for _ in range(n):
    c = random.random() < 0.01
    p = random.random() < (0.90 if c else 0.05)
    if p:
        pos += 1
        pos_c += c
print(pos_c / pos)   # 0.1521...   (calculation: 0.1538...)
\end{lstlisting}
\begin{itemize}
\item Agreement with $2/13$ checks that the \emph{code} implements the stated rates.
\end{itemize}
\note{Verified by running this exact script: 58,620 positives out of one million, and the share who actually have C is 0.1521. Hand calculation is 0.1538. Close, as expected from sampling noise. A second run with random.seed(10) on the coin flips gave 0.4953 for heads after five heads (62,137 qualifying flips) versus 0.5002 for heads overall; both near one half. Say what these checks are: the program and the table agree.}
\end{frame}

\begin{frame}{What a simulation does \emph{not} establish}
\begin{itemize}
\item It runs the assumptions you typed in. If 1\%, 90\%, and 95\% are wrong for real patients, the simulation is wrong in exactly the same way.
\item Agreement between table and simulation says nothing about whether the population matches the assumptions.
\item More runs shrink noise, not assumption error.
\end{itemize}
\vspace{2mm}
\begin{block}{Boundary note}
Calculation and simulation agree $\Rightarrow$ the code matches the model. They do \textbf{not} agree $\Rightarrow$ ``the model matches the world.''
\end{block}
\note{Thirteen minutes on this slide plus the previous one. Students write the boundary note: what was checked, what was not. This is the sentence the Decision Gate's Check section needs.}
\end{frame}

\begin{frame}{The AI base-rate check}
\begin{block}{AI answer to examine}
``A positive result is correct 90\% of the time, because the test is 90\% sensitive.''
\end{block}
\begin{itemize}
\item Which conditional probability did it confuse? ($P(+\mid C)$ for $P(C\mid +)$.)
\item What is missing from its Sources line? (The base rate.)
\item Independent check: the four-cell table gives $90/585\approx 15.4\%$.
\end{itemize}
\note{An AI proposal is not evidence until independently checked. The table is the independent check. Ask students to say which step an AI could have skipped silently, and how they would catch it.}
\end{frame}

\begin{frame}{The shape of every answer this week}
\begin{enumerate}
\item \textbf{Sources} --- rates, counts, rules, or simulation input.
\item \textbf{Rules / Assumptions} --- event definitions; independence, representativeness, randomness.
\item \textbf{Work} --- count table, calculation, or reproducible simulation.
\item \textbf{Check} --- totals, a second calculation, a changed base rate; say what it cannot validate.
\item \textbf{One-Sentence Summary} --- no stronger than the evidence.
\end{enumerate}
\vspace{1mm}
\small Example: ``In this 10{,}000-person example with 1\% prevalence, about 15 of every 100 positives have the condition; this depends on the stated rates holding for the population.''
\note{Draft the conclusion sentence now, as part of the 1:05 to 1:17 block. The example sentence reports the 1 percent result and the base-rate boundary, nothing more. Verified: 90/585 = 15.4 per 100 positives.}
\end{frame}

\begin{frame}{Decision Gate note, and Thursday}
\begin{itemize}
\item \textbf{Decision Gate:} one individual submission this week, using the five headings above. It is not an extra essay or a copy of Show \& Tell.
\item Choose the screening example, a stated game or drawing process, or an instructor-approved bounded claim.
\item Name where AI helped, and check the arithmetic and interpretation yourself.
\item \textbf{Thursday:} Show \& Tell defense with a calculation and an evidence table or simulation.
\end{itemize}
\small The due date and time are on the assignment in Canvas.
\note{Point to week-10/student/decision-gate.md and the reusable weekly problem-solving write-up. Today's trace (table, calculation, changed base rate, boundary note, summary sentence) is the start of the Gate. Do not state a deadline weekday here; the Canvas assignment is the source of truth.}
\end{frame}

\begin{frame}{Exit trace and what comes next}
\Large\centering
\textbf{Keep your trace: the 1\% table, $P(C\mid +)$, the 20\% prediction and result, your boundary note, and one bounded sentence.}
\normalsize\vspace{4mm}
\begin{itemize}
\item Next week: relations. Which pairs of things are related, and how do we classify the relation?
\end{itemize}
\note{1:17 to 1:25: inspect traces, collect questions. Bridge to Week 11: from probability over events to relations over pairs (reflexive, symmetric, transitive). Students keep the trace; it feeds the Decision Gate.}
\end{frame}

\end{document}
